\section{The Poisson hypertree}\label{sec:tree}

Seen from one agent, and with team deficits rescaled to $\cost_e=q^{2/3}\Defi(e)$, the cheap
teams of the market look like an infinite tree of teams in which every agent has infinitely
many teams but only finitely many below any cost level. This section defines that object,
the Poisson hypertree $\tree$, and the one notion of invariance used in this paper. It then
proves three things. Theorem~\ref{thm:decomp}: an invariant fractional perfect matching is
the marginal of an invariant random perfect matching of the same cost.
Lemma~\ref{lem:tree-duality}: weak duality between invariant fractional perfect matchings
and invariant prices. Theorem~\ref{thm:process}: a fixed point of the recursion map $\RDE$
gives an invariant perfect matching on $\tree$ whose cost is nine times the mean of the fixed
point. Existence of a fixed point, and a bound on its mean, are the subject of
Section~\ref{sec:existence}. The convergence of the market to $\tree$ is the subject of
Section~\ref{sec:localconv} and is not used here.

The construction follows the pattern of the recursive tree processes used for the random
assignment problem \cite{Aldous92,Aldous01,AldousSteele04}; invariance is the hypergraph
form of involution invariance \cite{AldousLyons07}. Two features need care: every agent has
infinite degree, and the matching is not known to be a function of the tree, so it has to be
built together with the tree by an extension theorem.

\emph{Standard tools.} The following theorems are used by name. From measure theory:
Kolmogorov's extension theorem for a sequence of standard Borel spaces, and the
Ionescu-Tulcea theorem \cite{Kallenberg21}. For Poisson processes: the
marking, mapping and restriction theorems; the fact that a Poisson process with diffuse
locally finite intensity is almost surely simple; and Mecke's formula for a Poisson process
with $\sigma$-finite intensity on a standard Borel space \cite{LastPenrose17}.

\subsection{Hypertrees and their codes}\label{subsec:tree-codes}

A \emph{hypertree} is a set of agents, each with a sector in $\mathbb{Z}/3$, and a set of
teams, each made of three agents, one from each sector, such that the incidence graph
(agents and teams as nodes, membership as edges) is a tree. Two agents $a,b$ are
\emph{adjacent}, written $a\sim b$, if they share a team; in a hypertree that team is unique
and is written $e(a,b)$. The \emph{distance} between two agents is the number of teams on the
path joining them. Every team $e$ carries a cost $\cost_e\in[0,\infty)$. The standing
condition is
\begin{itemize}
\item[(R0)] every agent is in infinitely many teams, their costs are pairwise distinct, and
only finitely many of them lie below any level.
\end{itemize}
An \emph{incidence} is a pair $(v,g)$ of an agent $v$ and a team $g\ni v$. Given an agent
$b\neq v$, the incidence $(v,g)$ \emph{points toward $b$} if $g$ is the first team on the
path from $v$ to $b$, and \emph{points away from $b$} otherwise; every incidence at $b$
itself points away from $b$. Given a team $e$, the incidence $(v,g)$ \emph{points toward $e$}
if $g=e$ or $g$ is the first team on the path from $v$ to the agents of $e$, and away from $e$
otherwise.

\emph{Rank coding.} Let $V=\bigcup_{n\ge0}(\N\times\{1,2\})^n$ be the set of addresses and
$|v|$ the length of $v\in V$. For $v\in V$, $k\ge1$ and $i\in\{1,2\}$ write $(v,k,i)$ for the
address $v$ followed by the letter $(k,i)$; for $v=\varnothing$ this is $(k,i)$. In a rooted
hypertree with (R0) the root $o$ gets the address $\varnothing$ and sector $s_o$. Its teams, in
increasing order of cost, are called $[\varnothing,1],[\varnothing,2],\dots$, and the two
other agents of $[\varnothing,k]$ get the addresses $(k,1)$ and $(k,2)$, where $(k,i)$ is the
one in sector $s_o+i$. Recursively, an agent $v\neq\varnothing$ has a \emph{parent team}
$e(v)$ (the first team on the path to $o$), and in it a \emph{parent agent} $p(v)$ (the agent
of $e(v)$ nearest to $o$) and a \emph{sibling} $\operatorname{sib}(v)$; its other teams, in
increasing order of cost, are $[v,1],[v,2],\dots$, and the other agents of $[v,k]$ are
$(v,k,1)$ and $(v,k,2)$, the $i$-th one in sector $s_v+i$. Write $\cost_{v,k}$ for the cost of
$[v,k]$. The \emph{code} of the rooted hypertree is
\[
  \omega=\bigl(s_o,(\cost_{v,k})_{v\in V,\,k\ge1}\bigr)\in\Omega:=\mathbb{Z}/3\times\mathcal{S}^V ,
\]
where $\mathcal{S}$ is the set of strictly increasing sequences in $[0,\infty)$ that tend to
infinity, a Borel subset of $[0,\infty)^{\N}$. Coding is a bijection between rooted hypertrees
with (R0), up to root-preserving isomorphism, and the Borel subset $\Omega_0\subset\Omega$ of
codes in which no child team of an agent has the same cost as the team through which the
agent was reached; the Poisson construction below takes values in $\Omega_0$ almost surely. The set $\Omega$ carries its
product Borel $\sigma$-field and is a standard Borel space.

A \emph{decoration} is a family of marks on teams, on agents or on incidences with values in a
Polish space (messages $X_{v\to g}$, weights $x_e$, indicators $m_e$, prices $u_v$). In the
code it is a family indexed by addresses, so a decorated code is a point of a countable product
of Polish spaces. A \emph{random decorated rooted hypertree} is a random decorated code.

\begin{remark}[truncated balls]\label{rem:tree-balls}
For $K>0$ let $T_K$ be the component of the root after the teams of cost above $K$ are deleted
(written $\tree_K$ for the Poisson hypertree), and let $B_r(T_K,o)$ be the rooted marked
hypergraph formed by the teams of $T_K$ that contain an agent at distance at most $r-1$ from
$o$ in $T_K$, with sectors and costs, up to isomorphism. The product $\sigma$-field on $\Omega$ equals the $\sigma$-field generated by the
maps $\omega\mapsto B_r(T_K,o)$, $r,K\in\N$. Indeed, a truncated ball is a measurable function
of the code: it reads coordinates selected by comparisons with $K$. Conversely, truncation at
$K$ removes the most expensive teams at each agent and so preserves ranks; hence
$\{\cost_{v,k}\le t\}$ is the event that for all large integers $K$ the ball $B_{|v|+1}(T_K,o)$
contains the agent with address $v$ and that agent has at least $k$ teams of cost at most $t$
other than its parent team.
\end{remark}

\subsection{Invariance}\label{subsec:tree-invariance}

Let $\mathfrak{n}$ be a decorated code and $b=(k,i)$ a neighbour of the root. Put
$\rho(\mathfrak{n},k,i)=(\mathfrak{n}',k',i')$, where $\mathfrak{n}'$ is the code of the same
decorated hypertree rooted at $b$, the marks staying on their teams, agents and incidences, and
$(k',i')$ is the address of $o$ in $\mathfrak{n}'$. Explicitly
$k'=1+\#\{j:\cost_{b,j}<\cost_{\varnothing,k}\}$ and $i'=3-i$, and every coordinate of
$\mathfrak{n}'$ is a coordinate of $\mathfrak{n}$ selected by finitely many such comparisons.
So $\rho$ is a measurable involution of the space of doubly rooted codes
$\{(\mathfrak{n},k,i)\}$. For a measurable function $f\ge0$ on that space write
\[
  f(\mathfrak{n};o,b)=f(\mathfrak{n},k,i),\qquad f(\mathfrak{n};b,o)=f(\rho(\mathfrak{n},k,i)).
\]

\begin{definition}[invariance]\label{def:tree-invariant}
A random decorated rooted hypertree $\mathfrak{N}$ is \emph{invariant} if for every measurable
$f\ge0$ on doubly rooted codes
\begin{equation}\label{eq:tree-MT}
  \E\sum_{b\sim o}f(\mathfrak{N};o,b)=\E\sum_{b\sim o}f(\mathfrak{N};b,o).
\end{equation}
The sums are over all neighbours $b=(k,i)$ of the root, and both sides may be $+\infty$.
\end{definition}

Identity \eqref{eq:tree-MT} is a mass-transport identity at the level of teams: mass is sent
only between agents that share a team. It is the only form of unimodularity used in this
paper. Its equivalence with mass transport between agents at arbitrary distance is not needed
and not claimed.

\begin{lemma}[null sets transfer]\label{lem:tree-null}
Let $\mathfrak{N}$ be invariant and let $A$ be a measurable set of decorated rooted codes with
$\P(\mathfrak{N}\in A)=0$. Then almost surely no re-rooting of $\mathfrak{N}$, at any agent,
lies in $A$.
\end{lemma}

\begin{proof}
Fix $K>0$ and let $f(\mathfrak{n};a,b)=\ind\{\cost_{e(a,b)}\le K\}\,\ind_A(\mathfrak{n}\text{
rooted at }a)$. By \eqref{eq:tree-MT},
\begin{multline*}
  \E\#\{b\sim o:\ \cost_{e(o,b)}\le K,\ \mathfrak{N}\text{ rooted at }b\text{ lies in }A\}\\
  =\E\bigl[\ind_A(\mathfrak{N})\cdot 2\,\#\{e\ni o:\cost_e\le K\}\bigr]=0 .
\end{multline*}
Let $K\to\infty$: almost surely no re-rooting at a neighbour of the root lies in $A$. So the
set $A_1$ of codes that have a re-rooting at a neighbour of the root in $A$ is measurable and
null. Apply the same argument to $A_1$ and induct on the distance from the root.
\end{proof}

\subsection{The Poisson hypertree and matchings on it}\label{subsec:tree-poisson}

Let $\Lam$ be the measure $\Lam(\dd\xi)=\tfrac12\xi^2\dd\xi$ on $[0,\infty)$.

\begin{definition}[Poisson hypertree]\label{def:tree-poisson}
Let $s_o$ be uniform on $\mathbb{Z}/3$ and let $(\Pi_v)_{v\in V}$ be independent Poisson
processes on $[0,\infty)$ with intensity $\Lam$, independent of $s_o$. Let $\cost_{v,k}$ be the
$k$-th point of $\Pi_v$. Since $\Lam([0,K])=K^3/6$ is finite for every $K$ and
$\Lam([0,\infty))=\infty$, the code $(s_o,(\cost_{v,k}))$ lies in $\Omega$ almost surely. The
random rooted hypertree with this code is the \emph{Poisson hypertree}
$\tree=(\tree,o,\cost)$.
\end{definition}

\begin{lemma}\label{lem:tree-poisson-invariant}
The Poisson hypertree is invariant.
\end{lemma}

The proof is given after Lemma~\ref{lem:tree-msg-invariant}, of which it is a special case.

Let $H$ be a hypertree. A \emph{fractional perfect matching} of $H$ is a family of weights
$x_e\in[0,1]$ on the teams with $\sum_{e\ni v}x_e=1$ at every agent $v$; it is a \emph{perfect
matching} if every $x_e$ is $0$ or $1$, that is, if the teams of weight one are pairwise
disjoint and cover every agent. An \emph{invariant fractional perfect matching} is an invariant
random decorated rooted hypertree $(T,o,\cost,x)$, possibly with further decorations, with
(R0), in which the decoration $x$ is almost surely a fractional perfect matching. It is
\emph{on $\tree$} if $(T,o,\cost)$ is the Poisson hypertree. By Lemma~\ref{lem:tree-null} it is
enough to require $x_e\in[0,1]$ for the teams at the root and $\sum_{e\ni o}x_e=1$ almost
surely. The decoration may use randomness beyond the tree and its costs. The \emph{cost} of
$x$ is
\begin{equation}\label{eq:tree-cost}
  c(x)=\E\sum_{e\ni o}\cost_e\,x_e\in[0,\infty] ,
\end{equation}
the expected cost of the root's team when $x$ is a perfect matching.
Remark~\ref{rem:tree-normalisation} relates $c$ to the deficit of an assignment in the market.

\subsection{Decomposition of fractional perfect matchings}\label{subsec:tree-decomp}

The deterministic part comes first. Let $H$ be a hypertree (countable degrees; (R0) and costs
are not needed), $x$ a fractional perfect matching of $H$ and $o$ an agent. Assign to every
agent $w$ a \emph{state} $t(w)$, which is a team containing $w$ or a cemetery state
$\dagger$, as follows.
\begin{itemize}
\item $t(o)=e$ with probability $x_e$, for $e\ni o$.
\item Let $w\ne o$ have parent team $f$ (the first team on the path to $o$) and parent agent
$p$. If $t(p)=f$ then $t(w)=f$. If $t(p)\ne f$ and $x_f<1$, then $t(w)=g$ with probability
$x_g/(1-x_f)$, for $g\ni w$, $g\ne f$; this is a probability vector because
$\sum_{g\ni w,\,g\ne f}x_g=1-x_f$. If $t(p)\ne f$ and $x_f=1$, then $t(w)=\dagger$.
\item Conditionally on the states of all agents at distance at most $n$ from $o$, the states
of the agents at distance $n+1$ are drawn independently.
\end{itemize}
So $t$ is a Markov chain indexed by the rooted tree; its law on the product $\sigma$-field
exists by the Ionescu-Tulcea theorem (one countable product of kernels per generation). Let
$m=\{t(w):w\text{ an agent},\ t(w)\ne\dagger\}$ and let $\mu^o_{H,x}$ be the law of $m$, a
probability measure on sets of teams.

\begin{lemma}[the chain]\label{lem:tree-chain}
\begin{enumerate}
\item[(i)] $\mu^o_{H,x}$-almost surely $m$ is a perfect matching of $H$.
\item[(ii)] $\mu^o_{H,x}(t(w)=g)=x_g$ for every incidence $(w,g)$; in particular
$\mu^o_{H,x}(e\in m)=x_e$ for every team $e$.
\item[(iii)] $\mu^o_{H,x}$ does not depend on $o$; write $\mu_{H,x}$.
\item[(iv)] If $\varphi:(H,x)\to(H',x')$ is an isomorphism, then
$\varphi_*\mu_{H,x}=\mu_{H',x'}$.
\end{enumerate}
\end{lemma}

\begin{proof}
\emph{Step 1 (consistency).} If $t(w)=g$ for one agent $w$ of a team $g$, then $t=g$ at all
three agents of $g$. Indeed, let $p$ be the agent of $g$ nearest to $o$; the other two agents
of $g$ have parent team $g$ and parent agent $p$. Such an agent has state $g$ only if
$t(p)=g$, and if $t(p)=g$ then both have state $g$. Hence two teams of $m$ that share an agent
$w$ both equal $t(w)$: the teams of $m$ are pairwise disjoint, and they cover every agent whose
state is not $\dagger$.

\emph{Step 2 (marginals).} We show $\P(t(w)=g)=x_g$ by induction on the distance from $o$. At
$o$ it is the definition. Let $w$ have parent team $f$ and parent agent $p$. First,
$\P(t(w)=f)=\P(t(p)=f)=x_f$. Let $g\ni w$, $g\ne f$. If $x_f<1$, then
$\P(t(w)=g)=\P(t(p)\ne f)\cdot x_g/(1-x_f)=x_g$. If $x_f=1$, then $x_g=0$, and $\P(t(w)=g)=0$
by the rule for that case. This proves (ii). Summing over $g\ni w$ gives
$\P(t(w)=\dagger)=0$ for every $w$, and there are countably many agents; with Step 1 this
proves (i). (The rule for $x_f=1$ is applied with probability $\P(t(p)\ne f)=0$, so any other
convention gives the same measure.)

\emph{Step 3 (root independence).} Let $f=\{o,o',o''\}$ be a team at $o$. If $x_f=1$ then under
$\mu^o$ and under $\mu^{o'}$ the three states are $(f,f,f)$ almost surely. Let $x_f<1$. Under
$\mu^o$ the triple $(t(o),t(o'),t(o''))$ equals $(f,f,f)$ with probability $x_f$ and equals
$(e,g,h)$, with $e,g,h$ all different from $f$, with probability $x_e\,x_g\,x_h/(1-x_f)^2$: the
root draws $e$ with probability $x_e$, and then $o'$ and $o''$ draw independently from their
kernels. This expression is symmetric in the roles of $o,o',o''$, so it is also the law of the
triple under $\mu^{o'}$. Given the triple, under either root the restrictions of $t$ to the
three branches hanging from $f$ (the branch at $a\in f$ consists of the agents separated from
$f$ by $a$, with $a$ itself) are conditionally independent, and the restriction to the branch
at $a$ is the chain started from $t(a)$ at $a$ and run with the transition rule above in which
``parent'' is read as ``neighbour toward $f$''. That rule is the same for the roots $o$ and
$o'$, because inside the branch at $a$ the directions toward $o$, toward $o'$ and toward $f$
coincide. So $\mu^o=\mu^{o'}$, and since $H$ is connected $\mu^o$ does not depend on $o$.

\emph{Step 4 (equivariance).} $\varphi$ carries the construction with root $o$ to the
construction with root $\varphi(o)$; now use (iii).
\end{proof}

\emph{Measurability.} In rank coding the construction is a measurable function
$m=\mathfrak{m}(\omega,x,U)$ of the code $\omega$, the weights $x$ and a family
$U=(U_v)_{v\in V}$ of independent uniform random variables, independent of $(\omega,x)$: the
agent at address $v$ uses $U_v$ and the cumulative sums of its kernel along the rank order of
its teams. (Where $x$ is not a fractional perfect matching put $m=\emptyset$.) For a finite set
of addresses and a prescribed pattern on it,
$(\omega,x)\mapsto\P(\mathfrak{m}(\omega,x,U)\text{ shows the pattern})$ is measurable by
Fubini's theorem, and such patterns generate the $\sigma$-field on sets of teams. So
$(\omega,x)\mapsto\mu_{(\omega,x)}$ is a probability kernel. By
Lemma~\ref{lem:tree-chain}(iii),(iv) it commutes with re-rooting: the image of
$\mu_{(\omega,x)}$ under the relabelling of addresses that $\rho$ induces is the measure $\mu$
at the re-rooted code.

\begin{theorem}[decomposition]\label{thm:decomp}
Let $(T,o,\cost,x)$ be an invariant fractional perfect matching. Conditionally on
$(T,o,\cost,x)$ draw $m$ from $\mu_{(T,\cost,x)}$, independently of everything else, and write
$m_e=\ind\{e\in m\}$. Then $(T,o,\cost,x,m)$ is invariant, $m$ is almost surely a perfect
matching, $\E[m_e\mid T,o,\cost,x]=x_e$ for every team $e$, and $c(m)=c(x)$ in $[0,\infty]$.
In particular every invariant fractional perfect matching on $\tree$ is the conditional
expectation, given the tree, the costs and $x$, of an invariant random perfect matching on
$\tree$ of the same cost.
\end{theorem}

\begin{proof}
Write $\mathfrak{n}$ for the code of $(T,o,\cost,x)$. \emph{Invariance.} Let $f\ge0$ be
measurable on doubly rooted codes decorated with $(x,m)$ and put
\[
  F(\mathfrak{n};a,b)=\int f\bigl((\mathfrak{n},m');a,b\bigr)\,\mu_{\mathfrak{n}}(\dd m').
\]
$F$ is measurable because $\mu$ is a kernel. Because the kernel commutes with re-rooting,
$\int f((\mathfrak{n},m');b,o)\,\mu_{\mathfrak{n}}(\dd m')$ is the value of $F$ at the
re-rooted code, that is, $F(\mathfrak{n};b,o)$ in the notation of
Section~\ref{subsec:tree-invariance}. Conditioning on $(T,o,\cost,x)$ and using Tonelli's
theorem,
\[
  \E\sum_{b\sim o}f(\cdot,m;o,b)=\E\sum_{b\sim o}F(\cdot;o,b),\qquad
  \E\sum_{b\sim o}f(\cdot,m;b,o)=\E\sum_{b\sim o}F(\cdot;b,o),
\]
and the two right-hand sides agree by \eqref{eq:tree-MT} for $(T,o,\cost,x)$.
\emph{Matching and marginals.} Lemma~\ref{lem:tree-chain}(i),(ii), applied conditionally on
$(T,o,\cost,x)$. \emph{Cost.} By Tonelli's theorem,
$c(m)=\E\sum_{e\ni o}\cost_e\,\E[m_e\mid T,o,\cost,x]=\E\sum_{e\ni o}\cost_e x_e=c(x)$.
\end{proof}

\subsection{Weak duality}\label{subsec:tree-duality}

\begin{lemma}[weak duality]\label{lem:tree-duality}
Let $(T,o,\cost,x,u)$ be \emph{jointly invariant}: one random decorated rooted hypertree with
(R0), with weights $x_e\in[0,1]$ on teams and prices $u_v\in\R$ on agents, that satisfies
\eqref{eq:tree-MT} for functions of all of $(T,\cost,x,u)$. Assume that almost surely
$\sum_{e\ni o}x_e=1$ and $u_o+u_b+u_c\le\cost_e$ for every team $e=\{o,b,c\}$ at the root, and
that $\E|u_o|<\infty$. Then
\[
  3\,\E u_o\le c(x)=\E\sum_{e\ni o}\cost_e\,x_e .
\]
\end{lemma}

\begin{proof}
Apply \eqref{eq:tree-MT} to $f_\pm(\mathfrak{n};a,b)=x_{e(a,b)}\,(u_b)_\pm$, where
$(\cdot)_\pm$ are the positive and negative parts:
\[
  \E\sum_{b\sim o}x_{e(o,b)}(u_b)_\pm=\E\sum_{b\sim o}x_{e(o,b)}(u_o)_\pm
  =\E\Bigl[(u_o)_\pm\cdot2\sum_{e\ni o}x_e\Bigr]=2\,\E(u_o)_\pm<\infty .
\]
So $\sum_{e=\{o,b,c\}\ni o}x_e(u_b+u_c)$ is integrable with mean $2\,\E u_o$, and, using
$\sum_{e\ni o}x_e=1$,
\[
  3\,\E u_o=\E\sum_{e=\{o,b,c\}\ni o}x_e\,(u_o+u_b+u_c)\le\E\sum_{e\ni o}x_e\,\cost_e .\qedhere
\]
\end{proof}

\begin{corollary}[prices that are a function of the tree]\label{cor:tree-factor}
Let $h$ be a measurable function of rooted codes of $(T,\cost)$ and put $u_v=h(T\text{ rooted
at }v)$. If $(T,o,\cost,x)$ is an invariant fractional perfect matching, the prices are almost
surely dual feasible at the root ($u_o+u_b+u_c\le\cost_e$ for every team $e=\{o,b,c\}$) and
$\E|u_o|<\infty$, then $3\,\E u_o\le c(x)$.
\end{corollary}

\begin{proof}
Every function of $(T,\cost,x,u;o,b)$ is a function of $(T,\cost,x;o,b)$, so $(T,o,\cost,x,u)$
is jointly invariant and Lemma~\ref{lem:tree-duality} applies.
\end{proof}

\begin{remark}\label{rem:tree-duality-joint}
Joint invariance is what the proof uses: the transported quantity $x_e u_b$ involves $x$ and
$u$ together. If $x$ and $u$ are each invariant together with $(T,\cost)$ but use separate
additional randomness with an unspecified coupling, \eqref{eq:tree-MT} for the pair is not
given.
\end{remark}

\subsection{The recursion map}\label{subsec:tree-rde}

Throughout the rest of this section $(\xi_k)_{k\ge1}$ are the points of a Poisson process on
$[0,\infty)$ with intensity $\Lam$, and $(X_k',X_k'')_{k\ge1}$ are independent pairs of
independent variables with law $\mu$, independent of $(\xi_k)$. For a probability law $\mu$ on
$\R$ such that $\inf_k(\xi_k-X_k'-X_k'')$ is almost surely finite, $\RDE(\mu)$ is the law of
this infimum. A \emph{fixed point} is a probability law $\mu$ on $\R$ for which the infimum is
almost surely finite and $\RDE(\mu)=\mu$. For a fixed point, $X,X',X'',X_1,X_2,X_3$ denote
independent variables with law $\mu$, and
\[
  Y=X'+X'',\qquad S=X_1+X_2+X_3 .
\]

\begin{lemma}[what the fixed-point equation forces]\label{lem:tree-G}
Let $\mu$ be a fixed point and put
\begin{equation}\label{eq:tree-G}
  G(x)=\tfrac16\,\E\bigl[(x+X'+X'')_+^3\bigr],\qquad x\in\R .
\end{equation}
\begin{enumerate}
\item[(i)] $G(x)<\infty$ for every $x$, and $\P(X>x)=\exp(-G(x))$. The function $G$ is convex
and twice continuously differentiable, with $G'(x)=\tfrac12\E(x+Y)_+^2>0$ and
$G''(x)=\E(x+Y)_+$, and it is an increasing bijection of $\R$ onto $(0,\infty)$. In particular
$\mu$ has no atoms.
\item[(ii)] Almost surely, for every $x$ only finitely many $k$ have $\xi_k-X_k'-X_k''\le x$,
the infimum is a minimum, and it is attained at exactly one $k$.
\item[(iii)] Fix $m$ with $p:=\P(Y>-m)>0$. Then $\P(X>x)\le\exp(-p(x-m)^3/6)$ for $x>m$, and
for every integer $n\ge1$
\[
  \E X_-^n\le\frac{\E Y_+^{n+3}}{(n+1)(n+2)(n+3)}<\infty .
\]
Hence $\E|X|^n<\infty$ for every $n$; in particular $\E|X|<\infty$ and $\E S_+^4<\infty$.
\end{enumerate}
\end{lemma}

\begin{proof}
Put $Y_k=X_k'+X_k''$. By the marking theorem the points $(\xi_k,Y_k)$ form a Poisson process
on $[0,\infty)\times\R$ with intensity $\Lam\otimes\mathrm{law}(Y)$. So
$N(x):=\#\{k:\xi_k-Y_k\le x\}$ is a Poisson variable with mean
\[
  \int_0^\infty\tfrac12\xi^2\,\P(Y\ge\xi-x)\dd\xi
  =\E\int_0^\infty\tfrac12\xi^2\,\ind\{\xi\le x+Y\}\dd\xi=G(x)\in[0,\infty] .
\]

(i) Suppose $G(x_0)=\infty$ for some $x_0$. For $x<x_0$ we have
$(x_0+Y)_+\le(x+Y)_++(x_0-x)$, so $(x_0+Y)_+^3\le4(x+Y)_+^3+4(x_0-x)^3$ and $G(x)=\infty$; and
$G$ is nondecreasing, so $G\equiv\infty$. A Poisson variable with infinite mean is almost
surely infinite, so $N(x)=\infty$ almost surely for every integer $x$ and the infimum is
$-\infty$ almost surely. This contradicts the hypothesis that the infimum is almost surely
finite. So $G$ is finite. Then $N(x+1)<\infty$ almost surely, so the infimum over the points
with $\xi_k-Y_k\le x+1$ is a minimum, and $\{\inf_k(\xi_k-Y_k)>x\}=\{N(x)=0\}$; hence
$\P(X>x)=\P(N(x)=0)=\exp(-G(x))$. Finiteness of $G$ gives $\E Y_+^3<\infty$, and
differentiation under the expectation (the derivatives $\tfrac12(x+Y)_+^2$ and $(x+Y)_+$ are
dominated, for $x\le x_0$, by their values at $x_0$) gives the formulas for $G'$ and $G''$ and
their continuity. $G$ is convex because $x\mapsto(x+y)_+^3$ is. Since
$\P(X>x)=\exp(-G(x))>0$ for every $x$, the variable $X$ is unbounded above; so is $Y$, and
therefore $\P(Y>-x)>0$ and $G'(x)>0$ for every $x$. Finally $G(x)\to0$ as $x\to-\infty$ by
dominated convergence and $G(x)\to\infty$ as $x\to\infty$ by monotone convergence.

(ii) $N(x)<\infty$ almost surely for every integer $x$, and $N$ is nondecreasing; so almost
surely finitely many points lie below every level and the infimum is attained. By the mapping
theorem the values $Z_k=\xi_k-Y_k$ form a Poisson process on $\R$ whose intensity measure has
distribution function $G$. As $G$ is continuous and finite, this measure is diffuse and locally
finite, and such a Poisson process is almost surely simple: no two $Z_k$ coincide. So the
minimiser is unique.

(iii) For $x>m$, $G(x)\ge\tfrac16\E[(x+Y)_+^3;\,Y>-m]\ge p(x-m)^3/6$. So $X_+$ has moments of
every order, and so has $Y_+\le X'_++X''_+$. For the lower tail,
$\P(X\le-t)=1-\exp(-G(-t))\le G(-t)=\tfrac16\E(Y-t)_+^3$, so
\begin{align*}
  \E X_-^n&=\int_0^\infty nt^{n-1}\,\P(X\le-t)\dd t
  \le\E\int_0^{Y_+}\frac{nt^{n-1}(Y-t)^3}{6}\dd t\\
  &=\frac{n\,\mathrm{B}(n,4)}{6}\,\E Y_+^{n+3}=\frac{\E Y_+^{n+3}}{(n+1)(n+2)(n+3)} ,
\end{align*}
where $\mathrm{B}(n,4)=6\,(n-1)!/(n+3)!$ is the Beta integral.
\end{proof}

So the tail condition that the construction below needs, $\E Y_+^3<\infty$, is a consequence
of $\RDE(\mu)=\mu$ and not an extra hypothesis. Equation \eqref{eq:tree-G}, read as an
equation for $G=-\ln\P(X>\cdot\,)$ in which $X',X''$ have survival function $\exp(-G)$, is
the fixed-point equation in the form used in Section~\ref{sec:existence}.

\subsection{From a fixed point to an invariant perfect matching}\label{subsec:tree-process}

\begin{theorem}[the tree process]\label{thm:process}
Let $\mu$ be a fixed point of $\RDE$. Then there is an invariant random perfect matching $M$
on the Poisson hypertree: a random decorated rooted hypertree $(\tree,o,\cost,M)$ that is
invariant, whose $(\tree,o,\cost)$-marginal is the Poisson hypertree, and in which $M$ is
almost surely a perfect matching. Its cost \eqref{eq:tree-cost} is
\[
  c(M)=\tfrac18\,\E S_+^4=9\,\E X ,\qquad\text{and moreover}\qquad \E S_+^3=6 .
\]
\end{theorem}

The proof occupies the rest of this subsection. The matching is defined from a family of
real \emph{messages}, one on each incidence, which solve the recursion \eqref{eq:tree-R}
below in both directions. The messages are constructed first on incidences that point toward
the root (Lemma~\ref{lem:tree-inward}), then on all incidences
(Lemma~\ref{lem:tree-messages}); invariance is proved through a Palm formula
(Lemmas~\ref{lem:tree-palm} and~\ref{lem:tree-msg-invariant}); the matching and its cost
follow. Throughout, $\mu$ is a fixed point and $G$ is the function \eqref{eq:tree-G}.

\subsubsection*{Inward messages}

\begin{lemma}[inward messages]\label{lem:tree-inward}
There is a probability space carrying the Poisson hypertree $(s_o,(\Pi_v)_{v\in V})$ and real
random variables $(X_v)_{v\in V}$ such that
\begin{enumerate}
\item[(a)] for every $n\ge0$ the variables $(X_v)_{|v|=n}$ are independent with law $\mu$ and
independent of $(s_o,(\Pi_w)_{|w|<n})$;
\item[(b)] almost surely, for every $v\in V$,
$X_v=\min_{k\ge1}\bigl(\cost_{v,k}-X_{(v,k,1)}-X_{(v,k,2)}\bigr)$, and the minimum is attained
at exactly one $k$.
\end{enumerate}
The joint law of $(s_o,(\Pi_v),(X_v))$ is uniquely determined by $\mu$.
\end{lemma}

For $v\ne\varnothing$, $X_v$ will be the message from $v$ into its parent team. $X_\varnothing$
is an auxiliary variable, the root's own minimum over all its teams.

\begin{proof}
\emph{Depth-$n$ laws.} For $n\ge0$ let $V_{<n}$ and $V_n$ be the sets of addresses of length
less than $n$ and equal to $n$, and
$\mathcal{E}_n=\mathbb{Z}/3\times\mathcal{S}^{V_{<n}}\times\R^{V_n}$. Both index sets are
countable, so $\mathcal{E}_n$ is a countable product of standard Borel spaces. Let $P_n$ be the
product measure on $\mathcal{E}_n$ under which $s_o$ is uniform, each $\Pi_w$ ($|w|<n$) is a
Poisson process of intensity $\Lam$ and each $X_v$ ($|v|=n$) has law $\mu$. A countable product
of probability measures needs no truncation of the cost level: this is where the infinite
degree is absorbed.

\emph{Bonding maps.} Define $\Phi_n:\mathcal{E}_{n+1}\to\mathcal{E}_n$ by keeping $s_o$ and
$\Pi_w$ for $|w|<n$ and setting, for $|v|=n$,
$X_v:=\inf_k(\cost_{v,k}-X_{(v,k,1)}-X_{(v,k,2)})$, with the value $0$ where the infimum is
$-\infty$. $\Phi_n$ is measurable.

\emph{Consistency: $\Phi_n(P_{n+1})=P_n$.} Under $P_{n+1}$ the blocks
\[
  \bigl(\Pi_v,(X_{(v,k,i)})_{k\ge1,\,i=1,2}\bigr),\qquad |v|=n,
\]
are independent of one another and of
$(s_o,(\Pi_w)_{|w|<n})$, because they are built from disjoint sets of coordinates of a product
measure. In each block $\Pi_v$ is a Poisson process of intensity $\Lam$ and the $X_{(v,k,i)}$
are independent with law $\mu$, independent of $\Pi_v$. So $X_v$, a function of the block of
$v$, has law $\RDE(\mu)=\mu$. Hence under $\Phi_n(P_{n+1})$ the variables $(X_v)_{|v|=n}$ are
independent with law $\mu$ and independent of $(s_o,(\Pi_w)_{|w|<n})$: this is $P_n$.

\emph{Extension.} For $N\ge0$ let $Q_N$ be the law on
$\mathcal{E}_0\times\dots\times\mathcal{E}_N$ of
$(\Phi_0\circ\dots\circ\Phi_{N-1}(\eta),\dots,\Phi_{N-1}(\eta),\eta)$ with $\eta\sim P_N$. By
consistency the projection of $Q_{N+1}$ on $\mathcal{E}_0\times\dots\times\mathcal{E}_N$ is
$Q_N$. Kolmogorov's extension theorem gives a probability measure $Q$ on
$\prod_n\mathcal{E}_n$ with these marginals. Under $Q$, almost surely
$\eta_n=\Phi_n(\eta_{n+1})$ for every $n$. So the coordinates $s_o$ and $\Pi_w$ of $\eta_n$ do
not depend on $n>|w|$. Let $X_v$ be the $X_v$-coordinate of $\eta_{|v|}$. Then (a) holds
because $\eta_n\sim P_n$, and the equation in (b) holds because the $X_v$-coordinate of
$\Phi_n(\eta_{n+1})$ is the infimum over the messages of level $n+1$. By (a) at level $|v|+1$
and Lemma~\ref{lem:tree-G}(ii), that infimum is almost surely a minimum attained at a unique
$k$; there are countably many $v$.

\emph{Uniqueness.} By (b), for every $n$ the family
$(s_o,(\Pi_w)_{|w|<n},(X_v)_{|v|\le n})$ is almost surely a fixed measurable function of
$(s_o,(\Pi_w)_{|w|<n},(X_v)_{|v|=n})$, whose law is $P_n$ by (a). These families generate the
$\sigma$-field.
\end{proof}

\subsubsection*{Branches}

A \emph{branch} is an agent with all its descendants, their costs and their inward messages.
Its code is $\beta=((\cost_{v,k})_{v\in V,\,k\ge1},(X_v)_{v\in V})$, without a sector, a point
of the standard Borel space $\mathfrak{S}:=\mathcal{S}^V\times\R^V$; $X(\beta):=X_\varnothing$
is its \emph{root message}. For $u\in V$ let $\mathfrak{X}_u$ be the branch at $u$ in the
process of Lemma~\ref{lem:tree-inward}, with addresses read relative to $u$, and let $L$ be the
law of $\mathfrak{X}_\varnothing$.

\begin{lemma}[independent branches]\label{lem:tree-branches}
For every $n\ge0$ the branches $(\mathfrak{X}_u)_{|u|=n}$ are independent with law $L$, and
independent of $(s_o,(\Pi_w)_{|w|<n})$. In particular $\mathfrak{X}_\varnothing$ consists of
$\Pi_\varnothing$ and of the branches $\mathfrak{X}_{(k,i)}$, which are independent with law
$L$ and independent of $\Pi_\varnothing$, together with
$X_\varnothing=\min_k\bigl(\cost_{\varnothing,k}-X(\mathfrak{X}_{(k,1)})-X(\mathfrak{X}_{(k,2)})\bigr)$.
\end{lemma}

\begin{proof}
Write $w\succeq u$ if $u$ is a prefix of $w$. Fix $N\ge n$. For $|u|=n$ the depth-$N$
truncation of $\mathfrak{X}_u$ and its input are
\[
  \bigl((\Pi_w)_{w\succeq u,\,|w|<N},(X_w)_{w\succeq u,\,|w|\le N}\bigr)
  \quad\text{and}\quad
  \bigl((\Pi_w)_{w\succeq u,\,|w|<N},(X_w)_{w\succeq u,\,|w|=N}\bigr).
\]
By Lemma~\ref{lem:tree-inward}(b) the truncation is one fixed measurable function of the
input. Under $P_N$ these inputs, for the
different $u$ of length $n$, are independent, independent of $(s_o,(\Pi_w)_{|w|<n})$, and each
has, after shifting addresses, the law of the corresponding input at the root under $P_{N-n}$.
So the truncations are independent, each with the law of the depth-$(N-n)$ truncation of
$\mathfrak{X}_\varnothing$, and independent of $(s_o,(\Pi_w)_{|w|<n})$. The truncations
increase with $N$ and generate the $\sigma$-field of the branch; a $\pi$-system argument
finishes the proof. The last sentence is the case $n=1$ with Lemma~\ref{lem:tree-inward}(b).
\end{proof}

Independence at fixed addresses is not the whole story for messages in both directions: the
$k$-th cheapest team at an agent is not independent of the other teams there, since those of
lower rank are cheaper. The correct two-directional statement is the Palm formula of
Lemma~\ref{lem:tree-palm}: seen from a typical team, its cost and the three branches hanging
from it are independent, and the branches are independent with law $L$.

\subsubsection*{Messages in both directions}

For a team $g$ and $v\in g$ write $g\setminus v$ for the set of the other two agents. The
target is a family $\mathfrak{M}=(X_{v\to g})$ of real numbers indexed by all incidences such
that, at every incidence $(v,g)$,
\begin{equation}\label{eq:tree-R}
  X_{v\to g}=\min_{g'\ni v,\ g'\neq g}\Bigl(\cost_{g'}-\sum_{w\in g'\setminus v}X_{w\to g'}\Bigr).
\end{equation}

\begin{definition}[the message family]\label{def:tree-messages}
\emph{Inward:} $X_{v\to e(v)}:=X_v$ for $v\ne\varnothing$. \emph{Outward,} by induction on
$|v|$:
\begin{align*}
  X_{\varnothing\to[\varnothing,k]}&:=\min_{j\ne k}\bigl(\cost_{\varnothing,j}-X_{(j,1)}-X_{(j,2)}\bigr),\\
  X_{v\to[v,k]}&:=\min\Bigl\{\cost_{e(v)}-X_{p(v)\to e(v)}-X_{\operatorname{sib}(v)},\
     \min_{j\ne k}\bigl(\cost_{v,j}-X_{(v,j,1)}-X_{(v,j,2)}\bigr)\Bigr\},
\end{align*}
the second line for $v\ne\varnothing$.
\end{definition}

By Lemma~\ref{lem:tree-inward}(a) and Lemma~\ref{lem:tree-G}(ii), almost surely at every agent
only finitely many of the values $\cost_{v,j}-X_{(v,j,1)}-X_{(v,j,2)}$ lie below any level. So
each of these minima is finite and attained, and each message is a measurable function of
$\mathfrak{X}_\varnothing$.

\begin{lemma}[the recursion in both directions]\label{lem:tree-messages}
\begin{enumerate}
\item[(i)] Almost surely, \eqref{eq:tree-R} holds for $\mathfrak{M}$ at every incidence.
\item[(ii)] Let $b$ be an agent. A family satisfying \eqref{eq:tree-R} at every incidence,
with all minima finite, is determined by its values on the incidences that point toward $b$.
The same holds with a team $e$ in place of $b$.
\end{enumerate}
\end{lemma}

\begin{proof}
(i) For the incidence $(v,e(v))$ the teams $g'\ne e(v)$ at $v$ are its child teams $[v,k]$,
whose other agents send their inward messages; so \eqref{eq:tree-R} at this incidence is
Lemma~\ref{lem:tree-inward}(b). For $(v,[v,k])$ the teams $g'\ne[v,k]$ at $v$ are the other
child teams and, if $v\ne\varnothing$, the parent team $e(v)$, into which $p(v)$ sends
$X_{p(v)\to e(v)}$ and $\operatorname{sib}(v)$ sends $X_{\operatorname{sib}(v)}$; so
\eqref{eq:tree-R} at this incidence is Definition~\ref{def:tree-messages}.

(ii) Let $(v,g)$ point away from $b$. In \eqref{eq:tree-R} the teams $g'\ne g$ at $v$ are of
two kinds. If $g'$ is not the first team on the path from $v$ to $b$, then $b$ and $v$ lie on
the same side of $g'$, so both incidences $(w,g')$, $w\in g'\setminus v$, point toward $b$. If
$g'$ is that first team (then $v\ne b$), $g'\setminus v$ consists of the next agent $p$ on the
path, with $(p,g')$ pointing away from $b$ and $p$ strictly closer to $b$ than $v$, and of a
third agent $s$, with $(s,g')$ pointing toward $b$. So $X_{v\to g}$ is a function of messages
that point toward $b$ and of at most one message that points away from $b$ at an agent closer
to $b$. Induct on the distance from $b$; at $v=b$ only the first kind occurs. For a team $e$
the argument is the same, by induction on the distance from $e$: at an agent $v\in e$ and an
incidence $(v,g)$ with $g\ne e$, the team $g'=e$ contributes the two messages into $e$ from the
other agents of $e$, which point toward $e$.
\end{proof}

So the family $\mathfrak{M}$ is an object on the unrooted hypertree: each message is the
stated function of the messages one step further out, in both directions on every incidence,
and re-rooting only relabels it. Let
\[
  \mathfrak{N}=(\tree,o,\cost,\mathfrak{M})
\]
be the Poisson hypertree decorated with all messages. It is a measurable function of
$(s_o,\mathfrak{X}_\varnothing)$.

\subsubsection*{The Palm formula and invariance}

For a branch code $\beta$ and $k\ge1$, $i\in\{1,2\}$, let $\beta^{(k,i)}$ be its sub-branch at
$(k,i)$, and let $\beta^{\setminus k}$ be the branch obtained by deleting the root's $k$-th
child team with its two sub-branches, re-indexing the remaining child teams by rank and
recomputing the root message as the minimum over the remaining teams. Both are measurable in
$\beta$. For the team $e=[\varnothing,k]$ at the root, with other agents $b=(k,1)$ and
$c=(k,2)$, the branches hanging from $e$ at $o,b,c$ are
$\mathfrak{X}_\varnothing^{\setminus k},\mathfrak{X}_b,\mathfrak{X}_c$, and their root messages
are $X_{o\to e},X_{b\to e},X_{c\to e}$ (Definition~\ref{def:tree-messages}, first line).

\begin{lemma}[Palm formula]\label{lem:tree-palm}
For every measurable $F\ge0$ on $[0,\infty)\times\mathfrak{S}^3$,
\begin{multline}\label{eq:tree-palm}
  \E\sum_{k\ge1}F\bigl(\cost_{\varnothing,k};\mathfrak{X}_\varnothing^{\setminus k},
     \mathfrak{X}_{(k,1)},\mathfrak{X}_{(k,2)}\bigr)\\
  =\int_0^\infty\Lam(\dd\xi)\int L(\dd\beta_0)L(\dd\beta_1)L(\dd\beta_2)\,
     F(\xi;\beta_0,\beta_1,\beta_2),
\end{multline}
and the same holds with $s_o$ as an additional argument of $F$ on the left and an independent
uniform element of $\mathbb{Z}/3$ on the right.
\end{lemma}

\begin{proof}
By Lemma~\ref{lem:tree-branches} and the marking theorem,
$\Xi:=\sum_k\delta_{(\cost_{\varnothing,k},\mathfrak{X}_{(k,1)},\mathfrak{X}_{(k,2)})}$ is a
Poisson process on $[0,\infty)\times\mathfrak{S}^2$ with intensity
$\Lam\otimes L\otimes L$, which is $\sigma$-finite. There is a measurable map $\mathfrak{b}$
from counting measures to branch codes with $\mathfrak{X}_\varnothing=\mathfrak{b}(\Xi)$ (rank
the points by their first coordinate, attach the marks, take the minimum as root message), and
then $\mathfrak{X}_\varnothing^{\setminus k}=\mathfrak{b}(\Xi-\delta_{x_k})$, where $x_k$ is
the $k$-th point of $\Xi$. Mecke's formula
\[
  \E\int\tilde F(x,\Xi-\delta_x)\,\Xi(\dd x)=\int\E\tilde F(x,\Xi)\,(\Lam\otimes L\otimes L)(\dd x),
\]
applied to $\tilde F((\xi,\beta_1,\beta_2),\varphi)=F(\xi;\mathfrak{b}(\varphi),\beta_1,\beta_2)$,
gives \eqref{eq:tree-palm}, because $\mathfrak{b}(\Xi)=\mathfrak{X}_\varnothing$ has law $L$.
The sector $s_o$ is independent of everything else.
\end{proof}

\begin{lemma}[invariance of the message process]\label{lem:tree-msg-invariant}
$\mathfrak{N}=(\tree,o,\cost,\mathfrak{M})$ is invariant in the sense of
Definition~\ref{def:tree-invariant}.
\end{lemma}

\begin{proof}
For an ordered triple $(a_1,a_2,a_3)$ of distinct agents forming a team $e$ define the
\emph{view}
\[
  \mathfrak{v}(a_1,a_2,a_3)=(\sigma;\cost_e;\beta_1,\beta_2,\beta_3),
\]
where $\sigma=(s_{a_1},s_{a_2},s_{a_3})$ is the triple of sectors, an ordering of
$\mathbb{Z}/3$, and $\beta_i$ is the branch hanging from $e$ at $a_i$: the agents separated
from $e$ by $a_i$, together with $a_i$, with their costs, coded by rank from $a_i$ outward, and
with the messages on their incidences that point toward $e$. The branch codes do not depend
on a root: inside the branch at $a_i$, ``toward $e$'' and ``toward $a_j$'' are the same
direction for every $j\neq i$.

\emph{Step 1: the doubly rooted network is a fixed function of the view.} There is one
measurable map $\Psi$ such that, almost surely, for every team $e$ at the root and every
ordering $(a_1,a_2,a_3)$ of its agents, $\Psi(\mathfrak{v}(a_1,a_2,a_3))$ is the code of
$\mathfrak{N}$ re-rooted at $a_1$ together with the address of $a_2$ in that code. Indeed, the
hypertree with its sectors and costs is obtained by gluing the three branches to a team of
cost $\cost_e$; the messages on incidences pointing toward $e$ are those carried by the
branches, the three messages into $e$ being the root messages $X(\beta_i)$; and the messages on
incidences pointing away from $e$ are obtained from them by the recursion \eqref{eq:tree-R},
outward from $e$, as in the proof of Lemma~\ref{lem:tree-messages}(ii). This defines $\Psi$
(with the value $0$ for a message whose minimum is not attained), and each coordinate of its
output is obtained from coordinates of the view by countably many comparisons and minima, so
$\Psi$ is measurable.
By Lemma~\ref{lem:tree-messages}(i) the family $\mathfrak{M}$ satisfies \eqref{eq:tree-R}
everywhere, so by Lemma~\ref{lem:tree-messages}(ii) the family reconstructed by $\Psi$ is
$\mathfrak{M}$, whichever agent of $e$ comes first.

\emph{Step 2: the view measure is symmetric.} For $o$ the root, $e=[\varnothing,k]$, $b=(k,i)$
and $c=(k,3-i)$ the third agent,
\[
  \mathfrak{v}(o,b,c)=\bigl((s_o,s_o+i,s_o-i);\cost_{\varnothing,k};
     \mathfrak{X}_\varnothing^{\setminus k},\mathfrak{X}_{(k,i)},\mathfrak{X}_{(k,3-i)}\bigr).
\]
Sum \eqref{eq:tree-palm} over $i\in\{1,2\}$, using that $s_o$ is uniform and independent and
that $L\otimes L$ is exchangeable. As $s_o$ and $i$ vary, $(s_o,s_o+i,s_o-i)$ runs once through
the six orderings of $\mathbb{Z}/3$. So for every measurable $h\ge0$
\[
  \E\sum_{b\sim o}h(\mathfrak{v}(o,b,c))
  =\frac13\sum_{\sigma}\int_0^\infty\Lam(\dd\xi)\int L^{\otimes3}(\dd\beta)\,
     h(\sigma;\xi;\beta_1,\beta_2,\beta_3),
\]
the sum being over the six orderings $\sigma$ of $\mathbb{Z}/3$. The right-hand side is
unchanged when $h$ is composed with the swap
\[
  \tau(\sigma;\xi;\beta_1,\beta_2,\beta_3)=\bigl((\sigma_2,\sigma_1,\sigma_3);\xi;\beta_2,\beta_1,\beta_3\bigr),
\]
because $L^{\otimes3}$ is exchangeable and the set of orderings is closed under the
transposition.

\emph{Step 3.} $\mathfrak{v}(b,o,c)=\tau(\mathfrak{v}(o,b,c))$. By Step 1,
$f(\mathfrak{N};o,b)=(f\circ\Psi)(\mathfrak{v}(o,b,c))$ and
$f(\mathfrak{N};b,o)=(f\circ\Psi)(\tau\,\mathfrak{v}(o,b,c))$. Apply Step 2 to $h=f\circ\Psi$
and to $h=f\circ\Psi\circ\tau$; the two results are equal, which is \eqref{eq:tree-MT}.
\end{proof}

The intensity $\Lam$ is $\sigma$-finite and not finite. That is harmless: \eqref{eq:tree-MT}
is an identity between elements of $[0,\infty]$.

\begin{proof}[Proof of Lemma~\ref{lem:tree-poisson-invariant}]
Repeat the proofs of Lemmas~\ref{lem:tree-branches}, \ref{lem:tree-palm}
and~\ref{lem:tree-msg-invariant} with the message coordinates deleted: a branch is then an
agent with its descendants and their costs, its law is that of
$(\Pi_v)_{v\in V}$, the branches at the agents of length $n$ are independent with that law by
construction, and Steps 1 to 3 need only the marking theorem and Mecke's formula. No law $\mu$
is involved.
\end{proof}

\subsubsection*{The matching}

\begin{definition}\label{def:tree-matching}
A team $e=\{a,b,c\}$ belongs to $M$ if and only if
$\cost_e<X_{a\to e}+X_{b\to e}+X_{c\to e}$.
\end{definition}

\begin{lemma}[the matching]\label{lem:tree-matching}
Almost surely $M$ is a perfect matching of $\tree$: every agent is in exactly one team of
$M$. The decorated hypertree $(\tree,o,\cost,M)$ is invariant and its
$(\tree,o,\cost)$-marginal is the Poisson hypertree. Moreover, almost surely, the team of $M$
at the root is the unique minimiser $e^*$ of
$e=\{o,b,c\}\mapsto\cost_e-X_{b\to e}-X_{c\to e}$ over the teams at the root, and the minimum
equals $X_\varnothing$.
\end{lemma}

\begin{proof}
For an agent $v$ and a team $g\ni v$ put $Z^v_g=\cost_g-\sum_{w\in g\setminus v}X_{w\to g}$. By
\eqref{eq:tree-R}, $X_{v\to e}=\min_{g\ni v,\,g\ne e}Z^v_g$. The defining inequality of $M$
reads $Z^v_e<X_{v\to e}$; so for each agent $v$ of a team $e$,
\begin{equation}\label{eq:tree-argmin}
  e\in M\iff Z^v_e<Z^v_g\ \text{ for every other team }g\text{ at }v .
\end{equation}
So at most one team of $M$ contains $v$, and exactly one does if and only if the minimum of
$(Z^v_g)_{g\ni v}$ is attained at a single team. At the root,
$Z^o_{[\varnothing,k]}=\cost_{\varnothing,k}-X_{(k,1)}-X_{(k,2)}$, and by
Lemma~\ref{lem:tree-inward}(b) the minimum over $k$ is almost surely attained at a single $k$
and equals $X_\varnothing$. The event that the minimum at the root is not attained at a single
team is therefore a null set of codes of $\mathfrak{N}$, and by
Lemmas~\ref{lem:tree-msg-invariant} and~\ref{lem:tree-null} it almost surely fails at every
agent. Hence almost surely every agent is in exactly one team of $M$.

$M$ is a measurable function of $\mathfrak{N}$ defined through incidences, with no reference
to the root. So every function of $(\tree,\cost,M;o,b)$ is a function of $(\mathfrak{N};o,b)$,
and \eqref{eq:tree-MT} for $(\tree,o,\cost,M)$ follows from
Lemma~\ref{lem:tree-msg-invariant}.
\end{proof}

\subsubsection*{Cost}

\begin{proof}[Proof of Theorem~\ref{thm:process}]
The matching $M$ of Definition~\ref{def:tree-matching} is an invariant random perfect matching
on $\tree$ by Lemma~\ref{lem:tree-matching}. It remains to compute its cost. Apply
\eqref{eq:tree-palm} to a function $F$ that depends on the branches only through their root
messages, which under $L^{\otimes3}$ are independent with law $\mu$: for measurable
$f\ge0$, and then for $f$ with $\int\Lam(\dd\xi)\,\E|f(\xi,X_1,X_2,X_3)|<\infty$ by
splitting $f$ into positive and negative parts,
\begin{equation}\label{eq:tree-campbell}
  \E\sum_{e=\{o,b,c\}\ni o}f(\cost_e,X_{o\to e},X_{b\to e},X_{c\to e})
  =\int_0^\infty\tfrac12\xi^2\,\E f(\xi,X_1,X_2,X_3)\dd\xi .
\end{equation}

\emph{First,} $f=\ind\{\xi<x_1+x_2+x_3\}$. The left-hand side of \eqref{eq:tree-campbell} is
$\E\#\{e\ni o:e\in M\}=1$, and the right-hand side is
$\E\int_0^{S_+}\tfrac12\xi^2\dd\xi=\E S_+^3/6$. So $\E S_+^3=6$.

\emph{Second,} $f=\xi\,\ind\{\xi<x_1+x_2+x_3\}$. The left-hand side is $c(M)$ and the right-hand
side is $\E\int_0^{S_+}\tfrac12\xi^3\dd\xi=\E S_+^4/8$, which is finite by
Lemma~\ref{lem:tree-G}(iii).

\emph{Third,} $f=(x_2+x_3)\,\ind\{\xi<x_1+x_2+x_3\}$. It is integrable:
$\int_0^\infty\tfrac12\xi^2\,\E[|X_2+X_3|\,\ind\{\xi<S\}]\dd\xi=\E[|X_2+X_3|\,S_+^3]/6$, and
$S_+^3\le9\sum_i(X_i)_+^3$, so this is at most a multiple of
$\E|X|\,\E X_+^3+\E|X|^4<\infty$. Hence the expected sum of the two partners' messages into the
root's team is $\E[(X_2+X_3)\,S_+^3]/6$. By exchangeability $\E[X_i\,S_+^3]$ does not depend on
$i$, so $\E[(X_2+X_3)\,S_+^3]=\tfrac23\E[S\,S_+^3]=\tfrac23\E S_+^4$, and the expectation equals
$\E S_+^4/9$.

\emph{Fourth.} Let $e^*=\{o,b^*,c^*\}$ be the root's team in $M$. By
Lemma~\ref{lem:tree-matching},
$X_\varnothing=\cost_{e^*}-X_{b^*\to e^*}-X_{c^*\to e^*}$. The variable $X_\varnothing$ has law
$\mu$ (Lemma~\ref{lem:tree-inward}(a) with $n=0$), and the terms on the right are integrable by
the second and third steps. So
\[
  \E X=\frac{\E S_+^4}{8}-\frac{\E S_+^4}{9}=\frac{\E S_+^4}{72},
  \qquad\text{that is,}\qquad c(M)=\frac{\E S_+^4}{8}=9\,\E X .\qedhere
\]
\end{proof}

\begin{remark}[a check of $\E S_+^3=6$ without the tree]\label{rem:tree-check}
Since $\P(X>x)=\exp(-G(x))$ and $G$ is continuous, $G(X)$ is a standard exponential variable.
With $X_3$ independent of $X_1+X_2$, which has the law of $Y$, equation \eqref{eq:tree-G}
gives $\E S_+^3/6=\E\,G(X_3)=1$.
\end{remark}

\begin{remark}[the matching is not claimed to be a function of the tree]\label{rem:tree-factor}
$M$ is a function of the tree, the costs and the messages. The messages are constructed
jointly with the tree by an extension theorem and are not claimed to be functions of
$(\tree,\cost)$; consequently $M$ is not claimed to be a function of $(\tree,\cost)$ either.
The definition of an invariant perfect matching allows such additional randomness, and
Theorem~\ref{thm:lp} is stated for matchings of this kind. The transfer to the market in
Section~\ref{sec:lp} works with conditional expectations of $M$ given the tree and the costs,
which are fractional; accordingly Theorem~\ref{thm:process} leads to a bound on the
fractional optimum $\LPq$ and not on $\Tq$. In fact the additional randomness cannot be
avoided: by Corollary~\ref{cor:proofmain-endogeny}, no invariant perfect matching on $\tree$
of cost below $\betaint$ is a measurable function of the tree and its costs.
\end{remark}

\begin{remark}[normalisation]\label{rem:tree-normalisation}
In the market a team has deficit $\Defi(e)$ and rescaled cost $\cost_e=q^{2/3}\Defi(e)$. Since
$\P(\Defi(e)\le t)=t^3/6$ for $t\le1$, an agent has on average
$q^2\cdot K^3/(6q^2)=K^3/6=\Lam([0,K])$ teams of rescaled cost at most $K$, for
$K\le q^{2/3}$; this fixes the intensity $\Lam$. A perfect assignment $N$ has $q$ teams and
every one of the $3q$ agents is in exactly one of them, so its total deficit is
\[
  \sum_{e\in N}\Defi(e)=q^{-2/3}\sum_{e\in N}\cost_e
  =q^{1/3}\cdot\frac{1}{3q}\sum_{\text{agents }v}\cost_{N(v)},
\]
where $N(v)$ is the team of $v$. So a total deficit of $c\,q^{1/3}$ means that $c$ is the
average cost per team, equivalently the average over agents of the cost of the agent's own
team. The tree quantity $c(m)=\E\sum_{e\ni o}\cost_e m_e$ is exactly this; the cost per agent
is $c/3$. The same holds for fractional $x$:
$\sum_e\Defi(e)x_e=q^{1/3}\cdot\frac{1}{3q}\sum_v\sum_{e\ni v}\cost_e x_e$, to be compared with
\eqref{eq:tree-cost}.
\end{remark}
